% ==============================================================================
% PLANTILLA MAESTRA GENERAL - UCNL
% ==============================================================================
% Uso:
% 1. Duplica este archivo y renombra la copia por materia, semana y actividad.
% 2. Actualiza documenttitle, documentsubtitle, documentsubject, coursename,
%    coursecode, universitydepartment y datos de la figura docente.
% 3. Lee la planeacion semanal y NO pegues las instrucciones completas dentro del
%    producto final. Usa la planeacion como lista de cumplimiento.
% 4. Revisa la bibliografia indicada, toma notas, registra citas y redacta con
%    comprension propia.
% 5. Identifica el producto academico solicitado y construyelo dentro del
%    documento siempre que sea posible.
% 6. Entrega en PDF con la nomenclatura indicada en el aula.
%
% Formato institucional recomendado:
% - Fuente tipo Helvetica, equivalente visual a Arial.
% - Interlineado 1.5.
% - Citas autor-anio con natbib: usar \citep{clave} o \citet{clave}.
% - Referencias en bibliografia BibTeX, formato APA 7 en cuerpo y referencias.
% - Caratula, resumen breve, desarrollo, producto solicitado, conclusion y
%   bibliografia.
% - Redaccion academica clara, coherente, cohesionada, formal y sin faltas.
% - La portada puede llevar una marca de agua institucional discreta; ajustar
%   las variables coverwatermark... si se requiere apagarla o cambiar intensidad.
%
% Criterios generales de cumplimiento:
% - Responder exactamente a la actividad solicitada en la planeacion.
% - Identificar conceptos, elementos, criterios o categorias de forma precisa.
% - Analizar, interpretar y tomar postura; no solo copiar fuentes.
% - Organizar la informacion por jerarquia, secuencia y criterios claros.
% - Incluir conclusion personal cuando la actividad lo permita o lo solicite.
% - Usar citas y referencias en APA 7.
% - Evitar copiar las indicaciones de la actividad en el cuerpo final.
%
% Regla para productos visuales:
% - Si la actividad pide tabla, cuadro comparativo, esquema, mapa conceptual,
%   mapa mental, linea de tiempo, diagrama, matriz, lista de verificacion, red
%   conceptual, SQA, organizador grafico o producto similar, realizarlo
%   preferentemente con LaTeX/TikZ para mantener nitidez, evidencia editable y
%   diseno controlado.
% - Si la consigna exige Canva, Miro u otra herramienta externa, insertar una
%   captura nitida y conservar la explicacion dentro del documento.
% - Para tablas extensas, usar pagina limpia, sin encabezados ni pies, landscape,
%   letra pequena y restaurar el estilo al terminar.
%
% Criterios de redaccion personal:
% - No escribir para llenar espacio; escribir para sostener una idea.
% - Convertir informacion en criterio: que significa, como funciona, que problema
%   resuelve, que limite presenta y que consecuencia produce.
% - Estructura mental recomendada: problema -> sistema -> consecuencia -> postura.
% - Usar pensamiento de diagnostico: entrada -> proceso -> salida.
% - Cada parrafo debe tener funcion: presentar, explicar, comparar, valorar o
%   concluir.
% - Usar primera persona con moderacion academica: "Considero que",
%   "Desde esta perspectiva", "La posicion que sostengo es". Evitar "yo creo".
%
% Notas de enfoque personal segun enfoque_personal.txt:
% - La voz debe ser academica, tecnica, directa, estrategica y funcional.
% - El texto no debe sonar como estudiante pasivo que recopila: debe sonar como
%   alguien que interpreta, ordena, cruza datos y busca consecuencia practica.
% - Antes de redactar, identificar la funcion de cada seccion: abrir problema,
%   definir criterio, demostrar con evidencia, comparar, valorar o concluir.
% - Evitar frases genericas como "es importante porque si" o "desde siempre".
%   Sustituirlas por relaciones verificables: causa, efecto, limite, costo,
%   funcion, riesgo, implicacion o mejora.
% - Cada concepto debe pasar de definicion a analisis. Formula util:
%   concepto -> funcion -> riesgo si se aplica mal -> consecuencia -> postura.
% - La postura personal no debe ser opinion suelta. Debe tener tres piezas:
%   posicion, razon y efecto. Ejemplo: "Considero que X, porque Y; por ello Z".
% - Integrar experiencia o criterio propio sin perder formalidad: no narrar la
%   vida personal, sino usar mirada practica para explicar por que el tema
%   importa en educacion, derecho, instituciones, comunicacion o profesion.
%
% Estilo observado en actividades ya realizadas:
% - Abrir con una tesis concreta:
%   "El Derecho no solo ordena conductas: organiza poder, define limites y
%   distribuye consecuencias."
% - Formular el problema en negativo y positivo:
%   "El problema no es solo..., sino..."
% - Convertir el producto visual en argumento:
%   "El cuadro no solo recopila datos; reconstruye la funcion, el problema y el
%   punto que requiere correccion."
% - Cerrar con aprendizaje transferible:
%   "No basta con saber que dice la norma; tambien debe analizarse a quien
%   protege, a quien excluye y que razones permiten defenderla publicamente."
%
% Nivel cognitivo esperado:
% - Recordar: usar solo para definir conceptos basicos. No debe ser el nivel
%   dominante en una actividad ponderable.
% - Comprender: explicar con palabras propias y relacionar con el tema.
% - Aplicar: llevar el concepto a un caso, norma, texto, correo, dilema o hecho.
% - Analizar: separar elementos, detectar causas, efectos, tensiones y limites.
% - Evaluar: sostener una postura con criterios, fuentes y consecuencias.
% - Crear: elaborar producto propio: mapa, cuadro, linea de tiempo, ensayo,
%   propuesta, reformulacion, glosario, portafolio o solucion argumentada.
%
% Regla de nivel minimo:
% - Si la actividad vale puntos y pide producto, no quedarse en recordar o
%   comprender. Debe llegar por lo menos a aplicar y analizar; si pide conclusion
%   o postura, debe llegar a evaluar.
%
% Uso de las 100 tecnicas didacticas:
% - Esta plantilla conserva el manual "Tecnicas didacticas del aprendizaje" que
%   se incorporo en las plantillas de Filosofia del Derecho, Etica y Moral
%   juridica, y Redaccion en contextos virtuales.
% - Cuando la planeacion mencione una tecnica de la coleccion UnADM, identificar:
%   definicion de la tecnica, producto esperado, pasos de construccion, criterios
%   de revision y cierre reflexivo.
% - La tecnica didactica no debe verse como adorno. Debe mostrar seleccion,
%   jerarquia, relacion entre ideas y lectura propia.
% - Siempre que sea viable, construir la tecnica en LaTeX/TikZ: cuadro
%   comparativo, mapa conceptual, esquema, mapa mental, linea de tiempo, matriz,
%   red conceptual, diagrama de flujo, tabla SQA o lista de verificacion.
%
% Checklist antes de entregar:
% [ ] El titulo, materia, semana, docente y fecha estan actualizados.
% [ ] El objetivo coincide con la planeacion.
% [ ] El producto solicitado aparece visible, rotulado y completo.
% [ ] La introduccion plantea un problema o postura, no solo anuncia el tema.
% [ ] El desarrollo explica, conecta y analiza.
% [ ] Hay citas APA autor-anio dentro del texto.
% [ ] La bibliografia final coincide con las fuentes citadas.
% [ ] La conclusion responde que se comprendio y que postura se sostiene.
% [ ] El nivel cognitivo llega al menos a aplicacion y analisis.
% [ ] Si hay conclusion/postura, incluye posicion, razon y consecuencia.
% [ ] La tecnica didactica solicitada esta construida y explicada.
% [ ] No aparecen instrucciones copiadas de la planeacion dentro del producto.
% [ ] Ortografia, acentos, sintaxis y nombres propios revisados.
% [ ] Si se uso IA como apoyo, el texto fue comprendido, revisado y apropiado.
% ==============================================================================

% CREACION DEL DOCUMENTO
\documentclass[
	spanish,
	letterpaper, oneside
]{article}

% ==============================================================================
% INFORMACION DEL DOCUMENTO
% Actualizar en cada actividad.
% Ejemplos:
% \def\documenttitle {Cuadro comparativo de conceptos fundamentales}
% \def\documentsubtitle {Actividad X - Administracion de emprendedores}
% \def\documentsubject {Informe academico UCNL}
% ==============================================================================
\def\documenttitle {Plantilla de reporte - Administracion de emprendedores}
\def\documentsubtitle {Actividad X - Administracion de emprendedores}
\def\documentsubject {Informe academico UCNL}

\def\documentauthor {Martin Jonathan de la Cruz}
\def\coursename {Administracion de emprendedores}
\def\coursecode {ADM-EMP}

\def\universityname {Universidad Ciudadana de Nuevo Leon}
\def\universityfaculty {Programa academico UCNL}
\def\universitydepartment {Administracion de emprendedores}
\def\universitydepartmentimage {departamentos/logo-ucnl}
\def\universitydepartmentimagecfg {height=1.57cm}
\def\universitylocation {Monterrey, Nuevo Leon}

% INTEGRANTES, PROFESORES Y FECHAS
\def\authortable {
	\begin{tabular}{ll}
		\textbf{Alumno:}
		& \begin{tabular}[t]{l}
			 Martin Jonathan de la Cruz \\
		\end{tabular} \\
		Matricula: & UCNL \\

		\textit{Figura docente:}
		& \begin{tabular}[t]{l}
			Nombre \\ Apellido
		\end{tabular} \\
		& \\
		\multicolumn{2}{l}{Fecha de realizacion: \today} \\
		\multicolumn{2}{l}{\universitylocation}
	\end{tabular}
}

% IMPORTACION DEL TEMPLATE BASE
\input{template}

% FORMATO GENERAL
\usepackage{helvet}
\renewcommand{\familydefault}{\sfdefault}

% TABLAS Y PRODUCTOS VISUALES
\usepackage{array}
\usepackage{longtable}
\usepackage{pdflscape}
\usepackage{tikz}
\usetikzlibrary{
	arrows.meta,
	positioning,
	calc,
	matrix,
	fit,
	shapes.geometric,
	shadows.blur
}

% CITAS EN FORMATO AUTOR-ANIO, COMPATIBLE CON APA 7 EN EL CUERPO DEL TEXTO
% Usar:
% - Cita parentetica: \citep{clave}
% - Cita narrativa: \citet{clave}
% - Varias fuentes: \citep{clave1,clave2}
\setcitestyle{authoryear,open={(},close={)}}

% ==============================================================================
% AYUDAS DE REDACCION
% Quitar o comentar \pendiente cuando el documento final este terminado.
% ==============================================================================
\newcommand{\pendiente}[1]{\textcolor{red}{[PENDIENTE: #1]}}

% ==============================================================================
% MARCA DE AGUA EN PORTADA
% - Se aplica solo a la portada porque usa \AddToShipoutPictureBG*.
% - Para apagarla en alguna actividad: \def\coverwatermarkenabled {false}
% - Usar opacidad baja para que no compita con titulo, datos ni logotipo.
% ==============================================================================
\def\coverwatermarkenabled {true}
\def\coverwatermarkimage {img/departamentos/logo-nuevo-leon.png}
\def\coverwatermarkopacity {0.16}
\def\coverwatermarkwidth {0.72\paperwidth}
\def\coverwatermarkxshift {0cm}
\def\coverwatermarkyshift {-0.35cm}

\newcommand{\insertcoverwatermark}{%
	\ifthenelse{\equal{\coverwatermarkenabled}{true}}{%
		\AddToShipoutPictureBG*{%
			\begin{tikzpicture}[remember picture,overlay]
				\node[
					opacity=\coverwatermarkopacity,
					inner sep=0pt,
					xshift=\coverwatermarkxshift,
					yshift=\coverwatermarkyshift
				] at (current page.center) {%
					\includegraphics[width=\coverwatermarkwidth]{\coverwatermarkimage}%
				};
			\end{tikzpicture}%
		}%
	}{}%
}

\begin{document}

% PORTADA
\insertcoverwatermark
\templatePortrait

% CONFIGURACION DE PAGINA Y ENCABEZADOS
\templatePagecfg
\onehalfspacing

% ==============================================================================
% RESUMEN / ABSTRACT
% Debe ser breve. Formula recomendada:
% Tema + objetivo + tecnica/producto + enfoque personal.
% ==============================================================================
\begin{abstractd}
	\textbf{Objetivo:} \pendiente{Redactar el objetivo con verbo academico: analizar, identificar, aplicar, comparar, explicar, evaluar, interpretar, argumentar o proponer. Debe coincidir con la planeacion semanal.}

	\textbf{Producto elaborado:} \pendiente{Indicar el producto solicitado: cuadro comparativo, mapa conceptual, linea de tiempo, estudio de caso, lista de verificacion, reporte, ensayo, foro, portafolio, glosario, resumen, esquema u otro.}

	\textbf{Enfoque de analisis:} \pendiente{Enunciar la postura central o criterio de lectura. Debe mostrar desde donde se analizara el tema y por que importa.}

	\textbf{Nivel cognitivo:} \pendiente{Indicar el nivel dominante de la actividad: comprender, aplicar, analizar, evaluar o crear. Evitar que el trabajo se quede solo en recordar informacion.}
\end{abstractd}

% TABLA DE CONTENIDOS - INDICE
\templateIndex

% CONFIGURACIONES FINALES
\templateFinalcfg

% ======================= INICIO DEL DOCUMENTO =======================

% ==============================================================================
% SECCION 1: INTRODUCCION
% Apertura recomendada:
% - Iniciar con una afirmacion clara, no con una frase generica.
% - Traducir la consigna a un problema academico, juridico, etico, comunicativo
%   o profesional, segun la materia.
% - Explicar por que importa.
% - Cerrar anticipando el enfoque personal.
% ==============================================================================
\section{Introduccion}

Emprender no consiste solo en tener una idea viable, sino en convertirla en una organizacion capaz de sostenerse bajo condiciones de incertidumbre, recursos limitados y presion competitiva. El problema central de esta actividad es que, en la practica, muchos proyectos se quedan en el nivel de entusiasmo inicial porque no traducen su propuesta en decisiones administrativas concretas: definir procesos, asignar recursos, medir resultados y ajustar estrategias con base en evidencia.

El objetivo de este reporte es analizar la administracion de emprendedores como un sistema de decision, no como una lista aislada de conceptos. Desde ese enfoque, se revisan los elementos que permiten pasar de una intuicion de negocio a un modelo operativo con capacidad de ejecucion y aprendizaje. La lectura del tema se orienta a identificar la funcion real de cada componente administrativo, sus limites cuando se aplica sin diagnostico y las consecuencias que produce en la viabilidad del emprendimiento.

La postura que guia el trabajo sostiene que la administracion emprendedora funciona cuando integra vision, estructura y control adaptativo; cuando alguno de estos elementos falta, el proyecto puede iniciar, pero dificilmente se consolida. Por ello, el analisis se centra en relacionar concepto, funcionamiento y efecto, con criterio tecnico y orientacion a resultados.

\section{Desarrollo}

\subsection{Contexto y problema}

\pendiente{Presentar el contexto de la actividad y transformar la consigna en problema de analisis.}

\subsection{Marco conceptual}

\pendiente{Explicar los conceptos centrales con fuentes. No acumular definiciones: cada concepto debe mostrar funcion, limite o consecuencia.}

\subsection{Nivel cognitivo aplicado}

\pendiente{Explicar brevemente que operacion cognitiva exige la actividad: comparar, diagnosticar, clasificar, interpretar, argumentar, evaluar o crear. Relacionar esa operacion con el producto solicitado.}

\subsection{Producto solicitado}

\pendiente{Insertar aqui el producto de la actividad. Si es visual, realizarlo en LaTeX/TikZ siempre que sea posible.}

\subsection{Tecnica didactica aplicada}

\pendiente{Cuando la planeacion indique una tecnica de las 100 Tecnicas Didacticas de Ensenanza y Aprendizaje, explicar aqui como se uso. Ejemplo: "La tecnica seleccionada organiza la informacion mediante criterios visibles, jerarquia y cierre reflexivo; por ello, el producto no solo presenta datos, sino que demuestra analisis y comprension" \citep{unadm100TecnicasDidacticas2023}.}

% ==============================================================================
% BLOQUE OPCIONAL: TABLA O CUADRO COMPARATIVO AMPLIO
% Usar cuando el producto necesite mucho ancho:
% - \clearpage para iniciar en pagina limpia.
% - \pagestyle{empty} para quitar encabezado y pie en el producto amplio.
% - landscape para aprovechar el ancho horizontal.
% - scriptsize, tabcolsep y arraystretch para controlar legibilidad.
% - Restaurar \pagestyle{fancy} al terminar.
%
% Si la tabla es corta, puede omitirse landscape y usarse longtable normal.
% ==============================================================================
% \clearpage
% \pagestyle{empty}
% \begin{landscape}
% \begingroup
% \scriptsize
% \setlength{\tabcolsep}{4pt}
% \renewcommand{\arraystretch}{1.35}
% \begin{longtable}{@{}p{0.18\linewidth}p{0.39\linewidth}p{0.39\linewidth}@{}}
% \caption{Titulo del cuadro comparativo}\label{tab:cuadro-comparativo}\\
% \toprule
% \textbf{Criterio} & \textbf{Elemento 1} & \textbf{Elemento 2} \\
% \midrule
% \endfirsthead
% \toprule
% \textbf{Criterio} & \textbf{Elemento 1} & \textbf{Elemento 2} \\
% \midrule
% \endhead
% Concepto & \pendiente{Contenido sintetico con cita si aplica.} & \pendiente{Contenido sintetico con cita si aplica.} \\
% Caracteristicas & \pendiente{Contenido.} & \pendiente{Contenido.} \\
% Funcion & \pendiente{Contenido.} & \pendiente{Contenido.} \\
% Valoracion critica & \pendiente{Contenido.} & \pendiente{Contenido.} \\
% \bottomrule
% \end{longtable}
% \endgroup
% \end{landscape}
% \pagestyle{fancy}

% ==============================================================================
% BLOQUE OPCIONAL: PRODUCTO VISUAL EN TIKZ
% Usar para mapa conceptual, mapa mental, esquema, flujo, proceso o diagrama.
% Cada flecha debe expresar relacion; no usar nodos sueltos sin conectores.
% ==============================================================================
% \begin{figure}[H]
% 	\centering
% 	\begin{tikzpicture}[
% 		node distance=1.3cm,
% 		box/.style={rectangle, rounded corners=2pt, draw=black!65, fill=black!4, align=center, minimum width=3.2cm, minimum height=0.9cm},
% 		arrow/.style={-{Latex[length=2.5mm]}, thick, draw=black!65}
% 	]
% 		\node[box] (centro) {Concepto central};
% 		\node[box, right=of centro] (relacion) {Concepto relacionado};
% 		\node[box, below=of relacion] (efecto) {Consecuencia};
%
% 		\draw[arrow] (centro) -- node[above]{implica} (relacion);
% 		\draw[arrow] (relacion) -- node[right]{produce} (efecto);
% 	\end{tikzpicture}
% 	\caption{Titulo del producto visual.}
% \end{figure}

\subsection{Analisis propio}

El producto elaborado permite observar un hallazgo clave: la diferencia entre un emprendimiento que avanza y otro que se estanca no suele estar en la idea original, sino en la calidad del sistema administrativo que la respalda. En terminos funcionales, planear define direccion, organizar asigna capacidades, dirigir alinea ejecucion y controlar corrige desviaciones; si estas funciones se aplican de forma fragmentada, la operacion pierde coherencia y la toma de decisiones se vuelve reactiva.

La interpretacion que sostengo es que la administracion de emprendedores debe leerse como una arquitectura de aprendizaje continuo. Su valor no esta solo en ordenar tareas, sino en reducir incertidumbre mediante ciclos de prueba, medicion y ajuste. Cuando el emprendedor decide con criterios verificables (costos, tiempos, respuesta del mercado, productividad), transforma la gestion en una ventaja competitiva; cuando decide por impulso, incrementa riesgo operativo, desgaste financiero y perdida de oportunidad.

La consecuencia practica de este analisis es directa: administrar bien no garantiza exito inmediato, pero administrar mal acelera el fracaso. Por eso, la postura tecnica que defiendo es que el emprendimiento requiere disciplina gerencial desde etapas tempranas, aun en proyectos pequenos. Esta disciplina no limita la innovacion; la hace sostenible, porque convierte la creatividad en procesos repetibles y en resultados evaluables.

\section{Postura personal}

Considero que la administracion de emprendedores debe asumirse como competencia estrategica y no como requisito academico accesorio. La razon es que emprender implica decidir bajo incertidumbre, y sin un marco administrativo las decisiones se vuelven inconsistentes, dificiles de evaluar y costosas de corregir.

Desde esta perspectiva, la formacion en esta materia aporta un criterio profesional transferible: antes de crecer, un proyecto debe demostrar que puede operar con orden, medir su desempeno y ajustar su rumbo sin perder su propuesta de valor. La consecuencia de adoptar esta postura es una practica emprendedora mas responsable, con mayor probabilidad de permanencia y con mejor uso de recursos humanos y financieros.

\section{Conclusion}

En este reporte se comprendio que la administracion de emprendedores no se limita a describir etapas del negocio, sino que articula un sistema para decidir, ejecutar y corregir con base en evidencia. El analisis permitio identificar que la funcion administrativa cumple un papel causal: orienta la accion, disminuye errores de implementacion y mejora la capacidad de adaptacion del proyecto frente a cambios del entorno.

Como cierre argumentativo, sostengo que un emprendimiento se vuelve viable cuando integra vision emprendedora con gestion tecnica. Esta posicion se fundamenta en que la improvisacion puede iniciar un proyecto, pero solo la administracion consistente permite sostenerlo, escalarlo y defender su permanencia. En consecuencia, la aportacion principal de la actividad fue pasar de una comprension descriptiva del tema a una postura evaluativa: emprender con metodo no es una opcion secundaria, es una condicion de supervivencia organizacional.

\bibliography{administracion-de-emprendedores}

% FIN DEL DOCUMENTO
\end{document}





